
\documentclass[11pt]{article}
\usepackage[margin=1in]{geometry}
\usepackage{amsmath,amssymb,amsthm,mathtools,bm}
\usepackage{physics}
\usepackage{hyperref}
\usepackage{enumitem}
\usepackage{microtype}
\usepackage{booktabs}

\newtheorem{theorem}{Theorem}[section]
\newtheorem{proposition}[theorem]{Proposition}
\newtheorem{lemma}[theorem]{Lemma}
\newtheorem{corollary}[theorem]{Corollary}
\newtheorem{conjecture}[theorem]{Conjecture}
\theoremstyle{remark}
\newtheorem{remark}[theorem]{Remark}

\newcommand{\Kep}{\mathrm{Kep}}
\newcommand{\Rminus}{R_-^2}
\newcommand{\Rplus}{R_+^2}
\newcommand{\1}{\mathbf 1}

\title{Regime Transitions on the Kepler Cone:\\
From the Audited Planar $R_-^2$ Sector to the 3D BPS Open EFT}
\author{Working Draft}
\date{\today}

\begin{document}
\maketitle

\begin{abstract}
We develop an audited open-system effective theory for a bound Kepler orbit
weakly coupled to a Maxwellian bath.  The analysis separates two stochastic
models that had previously been conflated.  In the planar constrained-noise
sector, with oriented angular momentum $\ell_z$, the corrected
half-covariance It\^o normalization gives a sign-changing secular
$R_-^2=2T(1-q)$ response and an extreme-eccentricity fold.  This sector is a
consistent planar model but is not the physical three-dimensional isotropic
plasma problem.

For an isotropic bath the physical angular momentum is
$L=|\bm r\times\bm v|$.  Its It\^o curvature changes the leading-log
generator, and exact hard collisions require the bounded jump-safe shape
coordinate
\[
s:=1-q=1-\frac{L\sqrt{-2h}}{\mu}.
\]
In the hot-bath limit $\Theta\to0$, the exact Maxwellian generator reduces
analytically to
\[
\left\langle\mathcal Ls\right\rangle
=
\frac{5D_0T}{2\mu q}(3q^2-1),
\]
so that
\[
q_0=\frac1{\sqrt3},
\qquad
e_0=\sqrt{\frac23}.
\]
This reproduces both the zero and the absolute drift coefficient of the
independent high-speed encounter result of Hamilton \& Modak
\cite{HamiltonModak2024}.  It provides a normalization benchmark for the
Rosenbluth half-covariance convention, the three-dimensional It\^o curvature,
the change of variables to $s$, and the Kepler time measure.  Retaining the
full Maxwellian speed dependence continues this known hot-limit zero into a
two-parameter drift-zero curve.  The three-dimensional leading-log shape fold is
\[
(q_c^{(0)},\Theta_c^{(0)},e_c^{(0)})
=
(0.59618619,\;0.32464127,\;0.80284620).
\]
We then match the long-distance Balescu--Lenard/BPS Fokker--Planck sector to
the exact short-distance Rutherford jump process.  For bound-preserving shape
evolution the matched endpoint is
\[
q_c^{\rm bound}
=
0.59618619+\frac{0.0111477}{L_{\rm C}}+O(L_{\rm C}^{-2}),
\qquad
\Theta_c^{\rm bound}
=
0.32464127+\frac{0.12838}{L_{\rm C}}+O(L_{\rm C}^{-2}).
\]
For the cemetery-extended killed observable, absorption produces the larger
endpoint displacement
\[
q_c^{\rm kill}
=
0.59618619+\frac{0.0683356}{L_{\rm C}}+O(L_{\rm C}^{-2}),
\qquad
\Theta_c^{\rm kill}
=
0.32464127-\frac{0.474482}{L_{\rm C}}+O(L_{\rm C}^{-2}).
\]
The latter is an unconditional killed-moment boundary, not the phase boundary
of the normalized surviving population.  The resulting EFT is therefore a
hierarchy: planar cone diffusion, three-dimensional jump-safe shape dynamics,
and BPS-matched finite jumps with a distinct absorbing channel.
\end{abstract}

\tableofcontents

\noindent\textbf{Main-result roadmap.}
The physical three-dimensional calculation has three checkpoints:
(i) the exact hot-bath match
$q_0=1/\sqrt3$ to Hamilton--Modak,
(ii) its finite-temperature continuation to the saddle-node
$(q_c,\Theta_c)=(0.59618619,0.32464127)$, and
(iii) the BPS exact-hard-jump correction to that fold.
The preceding planar cone calculation is retained as the normalization audit
and constrained-noise precursor.


\section{Scope and corrected status}

This note replaces the earlier secular-positivity draft.  The conservative
Kepler cone geometry, Levi--Civita regularization, It\^o projection from
$(h,\ell)$ to $(T,Y)$, and Kepler time measure survive the audit.  What changes
is the microscopic normalization of the transverse bath diffusion and all
downstream drift formulas that depended on it.

The corrected claims established here are:
\begin{enumerate}[leftmargin=2em]
\item the Maxwellian friction shape is unchanged;
\item the per-direction transverse It\^o diffusion eigenvalue carries an
additional factor $1/2$ relative to the old cone note;
\item Proposition~7.1 of \emph{Rest of the Cone} must be replaced as a whole;
\item pointwise positivity of $\dot R_-^2$ is false;
\item global secular positivity of $\overline{\dot R_-^2}$ is also false;
\item the secular sign boundary has a fold endpoint and two asymptotically
distinct branches.
\end{enumerate}

Sections 2--13 below first audit the equal-mass \emph{planar
constrained-noise} problem using the oriented angular-momentum component
$\ell=\ell_z$.  This is not yet the physical three-dimensional isotropic
plasma model developed in the later sections.  In the planar sector we choose
the $Y\ge0$ branch, so that
\[
q=Y/T=\sqrt{1-e^2}\ge0.
\]
The opposite orientation is related by $Y\mapsto-Y$ and exchanges the roles
of $R_-^2$ and $R_+^2$.  General mass ratio and full three-dimensional
$|\bm L|$ dynamics require separate reductions.

\section{Kepler cone and Levi--Civita amplitudes}

Let
\[
T=-\frac{\mu}{2h},
\qquad
Y=\frac{\ell}{\sqrt{-2h}},
\qquad
X^2=T^2-Y^2.
\]
Then
\[
e=\frac XT,
\qquad
q:=\sqrt{1-e^2}=\frac YT,
\qquad
T^2=X^2+Y^2.
\]

The nonnegative spinor amplitudes are
\[
\boxed{
R_\pm^2=2(T\pm Y).
}
\]
In particular,
\[
w:=R_-^2=2(T-Y)=2T(1-q).
\]

The quadratic structure of the cone will reappear near the
radial/collision edge $q\to0$ at fixed negative energy:
\[
1-e=\frac{q^2}{1+e}.
\]
This is the degenerate radial ellipse $e\to1$ with
$h=-\mu/(2T)<0$ fixed; it is not the zero-energy parabolic Kepler orbit.
and near apoapsis, in the half-angle coordinate
\[
u=\tan\frac{E-\pi}{2},
\]
one has
\[
z-z_{\min}=O(u^2).
\]

\section{Microscopic diffusion convention}

The velocity SDE is written
\[
d\bm v=\bm a(\bm v)\,dt+B\,dW,
\]
with
\[
\boxed{
D:=\frac12BB^{\mathsf T}.
}
\]
Hence
\[
d v_i\,d v_j=2D_{ij}\,dt.
\]

For an isotropic bath,
\[
D
=
D_\parallel\,\hat{\bm v}\hat{\bm v}^{\mathsf T}
+
D_{\perp,\mathrm{eig}}
(I-\hat{\bm v}\hat{\bm v}^{\mathsf T}),
\]
so
\[
\operatorname{Tr}D
=
D_\parallel+2D_{\perp,\mathrm{eig}}.
\]

The convention audit shows that the older bath notes mixed this half-covariance
tensor with the full Kramers--Moyal variance rate.  The corrected equal-mass
Maxwellian coefficients, in the Rosenbluth--MacDonald--Judd framework
\cite{Rosenbluth1957}, are
\begin{equation}
\boxed{
A(v)
=
2\mathcal K\frac{\Psi(\chi)}{v^2},
}
\end{equation}
\begin{equation}
\boxed{
D_\parallel(v)
=
\mathcal K\frac{G(\chi)}{v},
}
\end{equation}
\begin{equation}
\boxed{
D_{\perp,\mathrm{eig}}(v)
=
\frac{\mathcal K}{2v}
[\erf(\chi)-G(\chi)],
}
\end{equation}
where
\[
\chi=\frac{v}{\sqrt2\sigma},
\]
\[
\Psi(\chi)
=
\erf\chi-\frac{2\chi}{\sqrt\pi}e^{-\chi^2},
\]
\[
G(\chi)
=
\frac{\Psi(\chi)}{2\chi^2},
\]
and
\[
\mathcal K
=
\frac{4\pi\kappa_{\rm int}^2n_f\ln\Lambda}{m_t^2}.
\]

The project-local Rutherford encounter calculation gives
\[
\frac{d}{dt}
\langle|\Delta\bm v_\perp|^2\rangle
=
\frac{2\mathcal K}{u}
\]
in the fast-particle/cold-field limit.  Azimuthal symmetry splits this equally
between the two transverse covariance eigenvalues, so
\[
D_{\perp,\mathrm{eig}}
\to
\frac{\mathcal K}{2u},
\]
confirming the normalization above without importing external notation.

\section{General-phase It\^o generator}

For the oriented planar angular-momentum component,
\[
b_h=-vA+D_\parallel+2D_{\perp,\mathrm{eig}},
\]
\[
b_\ell=-A\frac{\ell}{v},
\]
and
\[
C_{hh}=2D_\parallel v^2,
\qquad
C_{h\ell}=2D_\parallel\ell.
\]

For
\[
T=-\frac{\mu}{2h},
\]
\[
T_h=\frac{2T^2}{\mu},
\qquad
T_{hh}=\frac{8T^3}{\mu^2},
\]
hence
\[
\boxed{
\mathcal LT
=
\frac{2T^2}{\mu}
(-vA+D_\parallel+2D_\perp)
+
\frac{8T^3v^2}{\mu^2}D_\parallel.
}
\]

For
\[
Y=\frac{\ell}{\sqrt{-2h}},
\]
\[
Y_h=\frac{YT}{\mu},
\quad
Y_\ell=\frac{Y}{\ell},
\quad
Y_{hh}=\frac{3YT^2}{\mu^2},
\quad
Y_{h\ell}=\frac{YT}{\mu\ell},
\quad
Y_{\ell\ell}=0,
\]
so
\[
\boxed{
\frac{\mathcal LY}{Y}
=
-A\left(\frac{Tv}{\mu}+\frac1v\right)
+
D_\parallel
\left(
\frac{3T}{\mu}
+
\frac{3T^2v^2}{\mu^2}
\right)
+
2D_{\perp,\mathrm{eig}}\frac{T}{\mu}.
}
\]

\section{Corrected cone Wilson functions}

Define
\[
z=\frac{Tv^2}{\mu},
\qquad
\Theta=\frac{\mu}{2T\sigma^2},
\qquad
\chi^2=\Theta z,
\]
and
\[
\Gamma
=
\mathcal K\frac{T^{3/2}}{\mu^{3/2}}.
\]

Then
\[
\mathcal LT=\Gamma T f_T,
\qquad
\mathcal LY=\Gamma Y f_Y,
\]
with
\begin{equation}
\boxed{
f_T
=
\frac{4\sqrt\Theta}{\sqrt\pi}e^{-\Theta z}
+
4\sqrt z\,(2-\Theta)G(\sqrt{\Theta z}).
}
\label{eq:fTcorr}
\end{equation}

and
\begin{equation}
\boxed{
f_Y
=
\frac1{\sqrt z}
\left[
\erf(\sqrt{\Theta z})
+
\{2+3z-4\Theta(z+1)\}
G(\sqrt{\Theta z})
\right].
}
\label{eq:fYcorr}
\end{equation}

Because
\[
R_-^2=2(T-Y),
\]
\[
\mathcal LR_-^2
=
2\Gamma T f_-,
\qquad
f_-=f_T-qf_Y.
\]
Thus
\begin{equation}
\boxed{
f_-
=
\frac{2(2-q)\sqrt\Theta}{\sqrt\pi}e^{-\Theta z}
+
\frac{G(\sqrt{\Theta z})}{\sqrt z}
\left[
z\{8-3q-2(2-q)\Theta\}
+
2q(2\Theta-1)
\right].
}
\label{eq:fminuscorr}
\end{equation}

\section{Direct $w(h,\ell)$ reconstruction}

A completely independent route begins from
\[
w(h,\ell)
=
-\frac{\mu}{h}
-\frac{\sqrt2\,\ell}{\sqrt{-h}}.
\]
Direct differentiation gives
\[
w_h=\frac{2T(2T-Y)}{\mu},
\qquad
w_\ell=-2\sqrt{\frac{T}{\mu}},
\]
\[
w_{hh}
=
\frac{T^2(16T-6Y)}{\mu^2},
\]
\[
w_{h\ell}
=
-\frac{2T^{3/2}}{\mu^{3/2}},
\qquad
w_{\ell\ell}=0.
\]

Substitution into the It\^o generator yields
\begin{equation}
\boxed{
\begin{aligned}
\frac{\mathcal Lw}{\mathcal K}
={}&
\frac{4T(2T-Y)}
{\sqrt\pi\,\mu v}
\chi e^{-\chi^2}
\\
&+
\Psi(\chi)
\left[
\frac{8T^3v}{\chi^2\mu^2}
-\frac{3T^2Yv}{\chi^2\mu^2}
-\frac{4T^2}{\mu v}
+\frac{2TY}{\mu v}
-\frac{2TY}{\chi^2\mu v}
+\frac{4Y}{v^3}
\right].
\end{aligned}
}
\label{eq:wcorrected}
\end{equation}

Equation \eqref{eq:wcorrected} reduces identically to
\eqref{eq:fminuscorr}.  This replaces Proposition~7.1 of the earlier cone
note.  The coefficient $4T(2T-Y)$ is retained, but the entire
$\Psi$ bracket changes.

\section{Kepler secular measure}

At leading adiabatic order,
\[
\varepsilon_{\rm bath}
=
\frac{P}{\tau_{\rm bath}}\ll1,
\]
so the osculating orbit is frozen over one Kepler period.

Using eccentric anomaly $E$,
\[
r=T(1-e\cos E),
\]
\[
z(E)
=
\frac{1+e\cos E}{1-e\cos E},
\]
and
\[
\boxed{
d\mu_e
=
\frac{1-e\cos E}{2\pi}\,dE.
}
\]

The exact identity
\[
z\,d\mu_e
=
\frac{1+e\cos E}{2\pi}\,dE
\]
implies
\[
\boxed{
\langle z\rangle_{\Kep}=1.
}
\]

The corrected secular coefficient is
\begin{equation}
\boxed{
F_-(q,\Theta)
=
\int f_-(z(E),q,\Theta)\,d\mu_e(E).
}
\label{eq:Fminus}
\end{equation}

\subsection{Analytic hot-bath check}

Using
\[
G(x)=\frac{2x}{3\sqrt\pi}+O(x^3),
\qquad
e^{-\Theta z}=1+O(\Theta),
\]
equation \eqref{eq:fminuscorr} gives
\[
f_-
=
\frac{2\sqrt\Theta}{3\sqrt\pi}
\left[
6-5q+(8-3q)z
\right]
+
O(\Theta^{3/2}).
\]
Since \(\langle z\rangle_{\Kep}=1\),
\begin{equation}
\boxed{
F_-(q,\Theta)
=
\frac{4(7-4q)}{3\sqrt\pi}\sqrt\Theta
+
O(\Theta^{3/2})
>0
}
\qquad
(\Theta\to0^+).
\label{eq:hotsecular}
\end{equation}
Thus secular growth in the sufficiently hot-bath regime is analytic, not
merely numerical.

\section{Pointwise sign transition}

At fixed physical phase,
\[
G(\sqrt{\Theta z})
\sim
\frac{1}{2\Theta z}
\]
as
\[
\Theta\to\infty.
\]
Equation \eqref{eq:fminuscorr} therefore gives
\begin{equation}
\boxed{
f_-
\longrightarrow
\frac{2q-(2-q)z}{z^{3/2}}.
}
\label{eq:pointwisecold}
\end{equation}

Hence sufficiently cold baths produce negative local drift whenever
\[
\boxed{
z>\frac{2q}{2-q}.
}
\]

The first physical orbit for which this can happen is obtained by setting
\[
z_p=\frac{1+e}{1-e}
=
\frac{2q}{2-q}.
\]
The physical solution is
\[
q_{\rm pt}
=
0.9559344860\ldots,
\]
or
\begin{equation}
\boxed{
e_{\rm pt}
=
0.2935800716\ldots.
}
\end{equation}

Thus local sign reversal begins already at moderate eccentricity, near
periapsis.

\section{Secular fold and negative band}

Numerical continuation of the secular zero set
\[
F_-(q,\Theta)=0
\]
finds a fold at
\begin{equation}
\boxed{
q_c
=
0.00929307\ldots,
}
\end{equation}
\begin{equation}
\boxed{
\Theta_c
=
54.2694\ldots,
}
\end{equation}
or
\begin{equation}
\boxed{
e_c
=
0.99995682\ldots.
}
\end{equation}

At the fold,
\[
F_-(q_c,\Theta_c)=0,
\qquad
F_\Theta(q_c,\Theta_c)=0,
\]
with
\[
F_q
=
127.527\ldots>0,
\]
and
\[
F_{\Theta\Theta}
=
4.17398\times10^{-4}>0.
\]

These nonzero numerical derivatives are the standard nondegeneracy
signature of a saddle-node/fold.  The numerical continuation finds no other
secular zero for
\[
q>q_c,
\]
over the scanned domain.  For
\[
0<q<q_c,
\]
two roots
\[
\Theta_-(q)<\Theta_+(q)
\]
appear, and
\begin{equation}
\boxed{
F_-<0
\qquad
\Theta_-(q)<\Theta<\Theta_+(q).
}
\label{eq:negativeband}
\end{equation}

Near the fold,
\[
\delta=q_c-q,
\]
and
\begin{equation}
\boxed{
\Theta_\pm
=
\Theta_c
\pm
781.701\sqrt{\delta}
+
9.530\times10^3\,\delta
+
O(\delta^{3/2}).
}
\label{eq:foldbranches}
\end{equation}

In eccentricity variables,
\[
\boxed{
\Theta_\pm(e)
=
\Theta_c
\pm
8108.71\sqrt{e-e_c}
+
O(e-e_c).
}
\]

\section{Lower asymptote: finite-temperature radial-edge boundary}

As
\[
q\to0
\]
at fixed finite $\Theta$, introduce
\[
y=\sqrt{\Theta z}.
\]
The limiting Kepler measure is
\[
d\mu_{0,\Theta}
=
\frac{4}{\pi\sqrt\Theta}
\frac{dy}{(1+y^2/\Theta)^2}.
\]

At $q=0$, the corrected local coefficient becomes
\[
f_-^{\rm par}
=
\frac{4\sqrt\Theta}{\sqrt\pi}e^{-y^2}
+
\frac{4(2-\Theta)y}{\sqrt\Theta}G(y).
\]

Hence
\begin{equation}
\boxed{
F_{\rm par}(\Theta)
=
\frac{16}{\pi}
\int_0^\infty
\frac{
\dfrac{e^{-y^2}}{\sqrt\pi}
+
\dfrac{2-\Theta}{\Theta}yG(y)
}{
(1+y^2/\Theta)^2
}\,dy.
}
\label{eq:Fpar}
\end{equation}

High-precision quadrature gives a positive zero at
\begin{equation}
\boxed{
\Theta_-^{(0)}
=
20.4956600\ldots.
}
\end{equation}

A scan over positive \(\Theta\) finds no additional radial-edge zero.  The
small-\(q\) numerical roots are consistent with the expansion
\begin{equation}
\boxed{
\Theta_-(q)
=
20.4956600
+
781.986\,q
+
o(q).
}
\label{eq:lowerbranch}
\end{equation}
The constant term is therefore an exact one-dimensional root problem, while
the displayed linear coefficient is presently a numerical asymptotic
coefficient rather than an analytic theorem.

Thus the entrance into the negative secular band is controlled by a
finite-temperature radial-edge instability.

\section{Upper asymptote: correlated cold/radial-edge boundary}

Set
\[
s=q\Theta,
\qquad
L=\ln\frac1q.
\]
For fixed $s>0$,
\begin{equation}
\boxed{
F_-(q,s/q)
=
-\frac4\pi L
+
H(s)
+
o(1).
}
\label{eq:fixeds}
\end{equation}

The negative logarithm comes from the overlap
\[
\sqrt q\ll u\ll1,
\]
where
\[
f_-\,d\mu
\sim
-\frac8\pi\frac{du}{u}.
\]

The leading negative logarithm in \eqref{eq:fixeds} follows analytically
from the overlap region.  Numerical extraction of the finite part, together
with the large-\(s\) matched form, gives
\begin{equation}
\boxed{
H(s)
=
\frac8\pi s
-
\frac4\pi\ln s
+
C
+
o(1),
}
\end{equation}
with
\[
C
=
3.86618\ldots.
\]
The coefficients of this large-\(s\) finite-part form have been checked
numerically over increasing \(s\), but a full written matched-asymptotic
derivation should be supplied before treating the constant term as a theorem.

The upper root therefore obeys
\[
2s-\ln s
=
L-C_*+o(1),
\qquad
C_*=\frac{\pi C}{4}=3.03649\ldots.
\]
Hence
\begin{equation}
\boxed{
s_+(q)
=
\frac12\ln\frac1q
+
\frac12\ln\ln\frac1q
-
1.86482
+
o(1),
}
\end{equation}
and
\begin{equation}
\boxed{
\Theta_+(q)
=
\frac1q
\left[
\frac12\ln\frac1q
+
\frac12\ln\ln\frac1q
-
1.86482
+
o(1)
\right].
}
\label{eq:upperbranch}
\end{equation}

Thus the upper boundary belongs to a distinct cold/radial-edge scaling sector.

\section{Complete corrected phase diagram}

The local and secular transitions are parametrically distinct.

Pointwise:
\[
e<0.29358007
\quad\Rightarrow\quad
\text{no asymptotic cold negative phase},
\]
while
\[
e>0.29358007
\quad\Rightarrow\quad
\text{some sufficiently fast phases have }f_-<0.
\]

Secular numerical continuation:
\[
e<0.99995682
\quad\Rightarrow\quad
F_->0
\quad\text{over the scanned domain},
\]
while
\[
e>0.99995682
\]
opens the negative band
\[
\Theta_-(e)<\Theta<\Theta_+(e).
\]

At extreme eccentricity,
\begin{equation}
\boxed{
20.49566+781.986\,q+o(q)
<
\Theta
<
\frac1q
\left[
\frac12\ln\frac1q
+
\frac12\ln\ln\frac1q
-
1.86482+\cdots
\right].
}
\end{equation}

The corrected secular phase structure is therefore
\[
\boxed{
\text{hot bath: }R_-^2\text{ grows},
}
\]
\[
\boxed{
\text{extreme-eccentricity intermediate/cold band: }R_-^2\text{ shrinks},
}
\]
\[
\boxed{
\text{ultra-cold at fixed }q>0:\ \text{numerically }R_-^2\text{ grows again}.
}
\]

\section{Numerical status and reproducibility}

The algebraic formulas through \eqref{eq:Fminus}, the pointwise cold limit,
and the hot-bath secular expansion are analytic.  The detailed secular phase
boundary is presently numerical/matched-asymptotic.

For ordinary eccentricities the integral \eqref{eq:Fminus} was evaluated in
eccentric anomaly.  Near the parabolic edge an independent half-angle
coordinate
\[
u=\tan\frac{E-\pi}{2}
\]
was used, with
\[
z=
\frac{(1-e)+(1+e)u^2}
{(1+e)+(1-e)u^2}
\]
and the exact symmetric time density
\[
d\mu_e
=
\frac{2}{\pi}
\frac{(1+e)+(1-e)u^2}{(1+u^2)^2}\,du.
\]
The two quadratures agree in their common region.

The fold was obtained from the simultaneous numerical equations
\[
F_-(q,\Theta)=0,
\qquad
\partial_\Theta F_-(q,\Theta)=0.
\]
The leading fold coefficient
\[
C_q=\sqrt{\frac{2F_q}{F_{\Theta\Theta}}}
\simeq781.701
\]
is stable under changes of finite-difference step and quadrature coordinate.
The next common branch-shift coefficient is less well conditioned numerically;
the quoted \(9.53\times10^3\) should be regarded as a working value, not a
validated constant.

Likewise, the parabolic threshold
\[
\Theta_-^{(0)}=20.4956600\ldots
\]
is the numerical root of the explicit one-dimensional integral
\eqref{eq:Fpar}.  The upper finite-part constant \(C\simeq3.86618\) is
numerically stable, but the complete written large-\(s\) matched derivation
remains to be supplied.

Accordingly, no digit in the phase-boundary constants should presently be
interpreted as an interval-certified error bound.

\section{EFT interpretation}

The natural reduced object is not a globally positive coupling but the
sign-changing secular response function
\[
F_-(q,\Theta).
\]
We use ``Wilson function'' below only in the provisional EFT sense: a full
Wilson-coefficient interpretation requires an explicit operator basis,
power counting, and matching prescription beyond leading logarithmic order.

Near the fold, the reduced normal form is
\[
F_-
\simeq
127.527\,(q-q_c)
+
2.08699\times10^{-4}
(\Theta-\Theta_c)^2.
\]
Thus the natural local scaling variable is
\[
\boxed{
\Xi
=
\frac{\Theta-\Theta_c}{\sqrt{q_c-q}}.
}
\]

The lower and upper boundaries are controlled by two different asymptotic
charts:
\[
\text{finite-}\Theta\text{ radial-edge chart}
\quad\longrightarrow\quad
\Theta_-^{(0)},
\]
and
\[
\text{correlated cold/radial-edge chart}
\quad\longrightarrow\quad
\Theta_+(q).
\]

This is a more genuine effective-theory structure than the discarded
monotonicity theorem.  Coarse graining does not produce a universal sign;
instead it reveals a regime transition that is invisible in the conservative
cone alone.

\section{Relation to previous work}
\label{sec:related-work}

The ingredients entering the present calculation have substantial precedents
in both stellar dynamics and plasma kinetic theory.  Our contribution is not
a new formulation of orbit-averaged relaxation, a new set of Rosenbluth
coefficients, or a new hard--soft matching scheme.  Rather, it is the
combination of these ingredients for a particular bounded Kepler shape
observable, together with the resulting finite-temperature drift topology and
its finite-hard-collision correction.  We summarize the closest points of
contact here.

\paragraph{Orbit-averaged relaxation in Kepler potentials.}
Two-dimensional Fokker--Planck descriptions in energy and angular momentum
are classical in stellar dynamics.  In particular, Cohn \& Kulsrud \cite{CohnKulsrud1978} developed
an orbit-averaged treatment in energy--angular-momentum space for stars around
a massive black hole, including angular-momentum diffusion and a loss-cone
boundary.  Subsequent non-resonant-relaxation calculations retain the same
basic structure: local velocity-space diffusion coefficients are transformed
to energy and angular momentum and then averaged along Keplerian orbits.
Modern work on galactic nuclei commonly uses the normalized angular momentum
\[
j=\frac{J}{J_c(E)}=\sqrt{1-e^2}
\]
as the eccentricity coordinate.  Consequently, neither orbit averaging in
$(E,J)$ nor the use of $J/J_c$ is new here.  The bounded variable
\[
s=1-j
\]
is chosen because it is convenient for the present drift-sign problem and,
unlike the unbounded cone amplitude $R_-^2$, remains admissible when exact
finite jumps approach the bound--unbound separatrix.

\paragraph{A known high-speed eccentricity-drift reversal.}
The closest direct predecessor to our leading-log drift calculation is the
recent wide-binary analysis of Hamilton \& Modak \cite{HamiltonModak2024}.  For weak, isotropic,
impulsive, penetrative encounters they derive a Fokker--Planck equation in
semimajor axis and squared eccentricity,
\[
\varepsilon=e^2,
\]
with first and second moments of the form
\[
A_\varepsilon=C(1-2\varepsilon),
\qquad
D_\varepsilon=2C\,\varepsilon(1-\varepsilon),
\]
for an encounter-dependent positive coefficient $C$.  Thus their
single-orbit eccentricity drift already reverses sign at
$\varepsilon=1/2$.  Applying It\^o's formula to
$s=1-\sqrt{1-\varepsilon}$ gives
\[
A_s
=
\frac{C}{4q}(3q^2-1),
\qquad
q=\sqrt{1-e^2},
\]
and therefore the corresponding zero
\[
q=\frac{1}{\sqrt3},
\qquad
e=\sqrt{\frac23}.
\]
Accordingly, an orbit-averaged eccentricity-drift reversal under isotropic
encounters is not itself a new result.

The distinction relevant here is the velocity regime.  Hamilton \& Modak
work in the impulsive high-field-star-speed limit in which the encounter
variance becomes independent of binary orbital phase.  The present
calculation retains the full speed dependence of the Maxwellian
Rosenbluth tensor.  In our variables this promotes the high-speed drift zero
to a two-parameter condition
\[
F_s(q,\Theta)=0.
\]
We find that the resulting zero set terminates at the finite-temperature
saddle-node
\[
(q_c,\Theta_c)
=
(0.59618619,\;0.32464127).
\]
To our knowledge, this finite-temperature continuation and its saddle-node
endpoint have not previously been reported.

\paragraph{Resonant and non-resonant relaxation.}
The galactic-nucleus literature also studies diffusion of
$j=\sqrt{1-e^2}$ by resonant and non-resonant relaxation.  These works are
important conceptual neighbors but should be distinguished from the present
model.  Resonant relaxation is driven by coherent torques between
orbit-averaged Keplerian wires, whereas the bath considered here produces
local two-body velocity scattering.  Non-resonant relaxation is closer:
modern calculations construct local velocity kicks, transform them to
$(E,J)$, and orbit-average them along the Keplerian wire.  Our leading-log
calculation belongs to this latter Chandrasekhar--Rosenbluth family, but
specializes to a homogeneous Maxwellian bath and asks a different question:
the sign topology of the first moment of the bounded shape $s$ as the ratio
of orbital to bath thermal speeds is varied.

\paragraph{Boltzmann descriptions and finite jumps.}
Finite-jump descriptions of perturbed Kepler orbits also have precedents.
For example, Collins \& Sari \cite{CollinsSari2008} formulate a Boltzmann equation for almost
circular binaries undergoing impulsive encounters and show that eccentricity
and inclination can undergo L\'evy-flight evolution.  Hence the use of a
Boltzmann collision operator, or the observation that finite encounters
contain information beyond a second-order Fokker--Planck truncation, is not
new.

Our hard-collision calculation differs in two respects.  First, the observable
is the bounded shape $s=1-J/J_c(E)$ over the full eccentricity range rather
than a near-circular eccentricity vector.  Second, the hard operator is not
introduced as an independent phenomenological model: it is used to replace
the short-distance Fokker--Planck surrogate inside a controlled weak-coupling
matching calculation.

\paragraph{Balescu--Lenard and BPS matching.}
Collective kinetic theories based on Balescu--Lenard equations
\cite{Lenard1960,Balescu1960} are well established in both plasma and stellar dynamics, including secular
Balescu--Lenard descriptions of quasi-Keplerian stellar systems
\cite{FouvryPichonMagorrian2017}.  These
self-gravitating theories describe resonant action-space diffusion generated
by the fluctuating stellar system itself and are physically distinct from
the homogeneous plasma bath considered here.

For the plasma problem, Brown, Preston \& Singleton (BPS) \cite{BPS2005} provide the
relevant controlled weak-coupling matching.  Their calculation combines the
short-distance Boltzmann sector with the long-distance collective sector and
determines finite next-to-leading-order transport coefficients without an
arbitrary collision cutoff.  BPS also construct the corresponding
Fokker--Planck description.  We therefore do not regard either the
hard--soft matching method or the BPS transport coefficients as new.

The new step in the present work is the observable-level projection of this
matched kinetic theory onto the Kepler shape drift.  Schematically, we use
\[
H_s^{\rm matched}
=
H_{\rm BPS\!-\!FP}^{\rm full}
+
\left(
H_{\rm hard}^{\rm exact}[s]
-
H_{\rm hard}^{\rm BPS\!-\!FP}[s]
\right),
\]
so that the Fokker--Planck representation of the hard sector is replaced,
rather than double counted, by the exact Rutherford jump generator acting on
$s$.  At the leading-log fold the finite reconstructed-FP contribution and
the exact hard residual have opposite signs.  The latter is large enough to
reverse the direction of the NLO fold displacement predicted by the
reconstructed Fokker--Planck operator alone.  We have not located a previous
BPS calculation projected onto a Kepler orbital-shape drift, or a previous
calculation of this hard-jump displacement of a finite-temperature Kepler
drift fold.

\paragraph{Scope of the novelty claim.}
The drift zero studied here is an observable-dependent It\^o quantity and
should not be interpreted as a coordinate-invariant thermodynamic phase
transition.  We therefore refer to it as a \emph{secular drift-sign fold} or
\emph{stochastic shape-regime boundary}.  The literature comparison above
also narrows the priority claim: eccentricity-drift reversal, normalized
angular-momentum variables, Kepler orbit averaging, finite-jump Boltzmann
dynamics, and BPS matching all have established precedents.  The results for
which we find no direct predecessor are (i) the continuation of the
high-speed isotropic-encounter drift zero through the full
velocity-dependent Maxwellian Rosenbluth problem to a finite-temperature
saddle-node, and (ii) the observable-level BPS hard-jump correction to that
fold, including the reversal of the NLO shift relative to the reconstructed
Fokker--Planck prediction.  We state these priority claims with the
qualification ``to our knowledge.''


The comparison above fixes the interpretation of the calculation that follows.
Before presenting the finite-temperature fold, we first establish an exact
analytic literature benchmark: the full three-dimensional Maxwellian
generator reduces in the hot-bath limit to the Hamilton--Modak drift not only
in its zero but in its absolute coefficient.  The finite-temperature
saddle-node and its observable-level BPS hard-jump displacement are then the
candidate new results.



\section{Main analytic benchmark: exact hot-bath match}
\label{sec:hot-bath-benchmark}

The closest literature benchmark admits an exact analytic match.  Consider
the hot-bath limit
\[
\Theta=\frac{\mu}{2T\sigma^2}\longrightarrow0
\]
at fixed Kepler orbit.  Then $\chi=v/(\sqrt2\sigma)\to0$ everywhere along
the orbit.  The Maxwellian Rosenbluth functions have the expansions
\[
\Psi(\chi)
=
\frac{4}{3\sqrt\pi}\chi^3+O(\chi^5),
\qquad
G(\chi)
=
\frac{2}{3\sqrt\pi}\chi+O(\chi^3).
\]
Consequently,
\[
D_\parallel
=
D_0+O(\sigma^{-3}v^2),
\qquad
D_{\perp,\mathrm{eig}}
=
D_0+O(\sigma^{-3}v^2),
\]
where
\[
\boxed{
D_0
=
\frac{2\mathcal K}{3\sqrt{2\pi}\,\sigma}.
}
\]
The friction is parametrically smaller,
\[
A(v)=O\!\left(\frac{\mathcal K v}{\sigma^3}\right),
\]
so at leading order the local velocity process is simply isotropic additive
diffusion,
\[
d v_i\,d v_j=2D_0\delta_{ij}\,dt.
\]
This is precisely the phase-independent kick limit relevant to the
penetrative high-speed encounter calculation of Hamilton \& Modak
\cite{HamiltonModak2024}.

The match can be performed directly in the present variables.  Let
\[
h=\frac{v^2}{2}-\frac{\mu}{r},
\qquad
L=|\bm r\times\bm v|,
\qquad
q=\frac{L\sqrt{-2h}}{\mu},
\qquad
s=1-q.
\]
For constant isotropic diffusion $D_0$, the local It\^o coefficients are
\[
\mathcal L h=3D_0,
\qquad
\mathcal L L=D_0\frac{r^2}{L},
\]
and
\[
C_{hh}=2D_0v^2,
\qquad
C_{hL}=2D_0L.
\]
The required derivatives of $s(h,L)$ are
\[
s_h=-\frac{q}{2h},
\qquad
s_L=-\frac{q}{L},
\qquad
s_{hh}=\frac{q}{4h^2},
\qquad
s_{hL}=-\frac{q}{2hL},
\qquad
s_{LL}=0.
\]
Therefore
\[
\mathcal Ls
=
qD_0
\left[
-\frac{5}{2h}
-\frac{r^2}{L^2}
+\frac{v^2}{4h^2}
\right].
\]

For a Kepler ellipse of semimajor axis $T$,
\[
h=-\frac{\mu}{2T},
\qquad
L^2=\mu Tq^2,
\]
and the exact time averages are
\[
\langle v^2\rangle=\frac{\mu}{T},
\qquad
\langle r^2\rangle
=
T^2\left(1+\frac32e^2\right)
=
\frac{T^2}{2}(5-3q^2).
\]
It follows immediately that
\[
\boxed{
\left\langle\mathcal Ls\right\rangle
=
\frac{5D_0T}{2\mu q}(3q^2-1).
}
\]
Thus the hot-bath zero is
\[
\boxed{
q_0=\frac1{\sqrt3},
\qquad
e_0=\sqrt{\frac23}.
}
\]

This agrees exactly with the independent wide-binary result.  Hamilton \&
Modak write $\varepsilon=e^2$ and obtain
\cite{HamiltonModak2024}
\[
A_\varepsilon=C(1-2\varepsilon),
\qquad
D_\varepsilon=2C\varepsilon(1-\varepsilon).
\]
Applying It\^o's formula to
$s=1-\sqrt{1-\varepsilon}$ gives
\[
A_s=\frac{C}{4q}(3q^2-1).
\]
In our isotropic hot limit,
\[
\frac{\langle|\Delta\bm v|^2\rangle}{\Delta t}=6D_0.
\]
Their coefficient
\[
C
=
\frac{5T}{3\mu}
\frac{\langle|\Delta\bm v|^2\rangle}{\Delta t}
\]
therefore becomes
\[
C=\frac{10D_0T}{\mu},
\]
and hence
\[
\frac{C}{4q}(3q^2-1)
=
\frac{5D_0T}{2\mu q}(3q^2-1),
\]
identically equal to the result above.

This is a nontrivial normalization benchmark: the Rosenbluth
half-covariance convention, the three-dimensional It\^o curvature of $L$,
the transformation to $s=1-J/J_c(E)$, and the Kepler time measure all combine
to reproduce an independently derived encounter-theory result, including its
absolute coefficient once the kick-variance conventions are identified.

The finite-temperature calculation can therefore be viewed as the controlled
continuation
\[
q_0=\frac1{\sqrt3}
\quad\longrightarrow\quad
F_s(q,\Theta)=0,
\]
whose zero set terminates at
\[
(q_c,\Theta_c)=(0.59618619,0.32464127).
\]


\section{Three-dimensional jump-safe completion}

The preceding sections establish the audited \emph{planar} $R_-^2$ sector.
That sector remains a valid constrained-noise model, but it is not by itself
the physical BPS plasma problem: an isotropic three-dimensional collision
changes the orbital plane.  The exact jump theory therefore requires
\[
L=|\bm r\times\bm v|
\]
rather than an oriented planar component $\ell$.

The globally admissible shape coordinate is
\[
\boxed{
s:=1-q
=
1-\frac{L\sqrt{-2h}}{\mu},
\qquad 0\le s\le1.
}
\]
It is related to the planar cone amplitude by $R_-^2=2Ts$ before plane
tilting, but unlike $R_-^2$ it remains bounded at the Kepler separatrix
$h\to0^-$.  This distinction matters only when finite hard jumps are retained:
the exact hard generator of $R_-^2$ diverges logarithmically from the bound
side of $h'=0$, whereas the exact hard generator of $s$ is finite.

For the three-dimensional isotropic diffusion tensor the angular-momentum
magnitude has the It\^o curvature
\[
\frac{\partial^2 L}{\partial v_z^2}=\frac{r^2}{L}
\]
at a planar state, and hence
\[
\mathcal L L
=
-\frac{A L}{v}
+
D_\perp\frac{r^2}{L}.
\]
This term is absent for the linear oriented component $\ell_z$.  It changes
the leading-log secular theory before any NLO matching is performed.

The resulting jump-safe leading-log fold is
\[
\boxed{
q_c^{(0)}=0.59618619,
\qquad
\Theta_c^{(0)}=0.32464127,
\qquad
e_c^{(0)}=0.80284620.
}
\]
The local fold derivatives are
\[
F_q=2.89734555,
\qquad
F_{q\Theta}=3.80451347,
\qquad
F_{\Theta\Theta}=0.58826836.
\]

Thus the extreme-eccentricity planar $R_-^2$ fold derived above and the
three-dimensional $s$ fold are not competing numerical estimates of the same
quantity.  They belong to different stochastic models and different
observable domains.

\section{Full BPS matching for the bound-survivor shape fold}

Write the secular generator as
\[
\dot s
=
\mathcal K_0
\left[
L_{\rm C}F_s(q,\Theta)
+
H_s^{\rm BPS}(q,\Theta)
+
O(L_{\rm C}^{-1})
\right],
\]
where $L_{\rm C}$ denotes the large Coulomb logarithm.

At the three-dimensional LL fold the finite BPS reconstructed-FP term is
\[
H_{\rm FP}=+0.04519133,
\qquad
\partial_\Theta H_{\rm FP}=-0.02890912.
\]
Replacing the hard FP sector by the exact importance-sampled Rutherford jump
gives, for the bound-preserving process,
\[
H_{\rm hard}=-0.07750043,
\qquad
\partial_\Theta H_{\rm hard}=-0.08808845.
\]
Here $H_{\rm hard}$ is a \emph{residual},
\[
H_{\rm hard}
=
H_{\rm hard}^{\rm exact,bound}
-
H_{\rm hard}^{\rm BPS\!-\!FP},
\]
not an additional hard contribution.  Thus the bookkeeping is
\[
\boxed{
H_{\rm FP}^{\rm full}+H_{\rm hard}
=
H_{\rm soft}^{\rm BPS\!-\!FP}
+
H_{\rm hard}^{\rm exact,bound},
}
\]
so the hard sector is replaced rather than double counted.  In the
importance-sampled implementation, escaping hard events have zero exact
bound-preserving increment but their local FP surrogate is still subtracted;
this is precisely the subtraction required by the equation above.

The quoted hard derivative is numerical.  Its sampling uncertainty dominates
the uncertainty of the $\Theta$-shift and should not be interpreted as an
interval-certified value.

Therefore
\[
\boxed{
H_{\rm BPS}^{\rm bound}
=
-0.03230910,
}
\qquad
\boxed{
\partial_\Theta H_{\rm BPS}^{\rm bound}
=
-0.11699757.
}
\]

Let
\[
q_c=q_c^{(0)}+\frac{q_c^{(1)}}{L_{\rm C}}+\cdots,
\qquad
\Theta_c=\Theta_c^{(0)}+\frac{\Theta_c^{(1)}}{L_{\rm C}}+\cdots.
\]
The fold equations imply
\[
q_c^{(1)}=-\frac{H_c}{F_q},
\qquad
\Theta_c^{(1)}
=
-\frac{F_{q\Theta}q_c^{(1)}+H_\Theta}{F_{\Theta\Theta}}.
\]
Hence the full bound-survivor BPS fold is
\[
\boxed{
q_c^{\rm bound}(L_{\rm C})
=
0.59618619+\frac{0.0111477}{L_{\rm C}}+O(L_{\rm C}^{-2}),
}
\]
\[
\boxed{
\Theta_c^{\rm bound}(L_{\rm C})
=
0.32464127+\frac{0.12838}{L_{\rm C}}+O(L_{\rm C}^{-2}),
}
\]
and
\[
\boxed{
e_c^{\rm bound}(L_{\rm C})
=
0.80284620-\frac{0.0082782}{L_{\rm C}}+O(L_{\rm C}^{-2}).
}
\]

The sign is nontrivial.  The reconstructed-FP finite term alone shifts the
fold toward smaller $q$, but the exact finite hard jumps overcompensate that
motion and leave a smaller net shift toward larger $q$.

\section{Cemetery-weighted killed-moment endpoint}

For an absorbing population, an unbound hard jump is assigned the cemetery
value $s'=0$.  At the LL $s$ fold the exact hard escape hazard is
\[
\frac{\lambda_{\rm esc}}{\mathcal K_0}
=
0.41027\pm0.00088
\]
in the present importance-sampled normalization.  Because
\[
s_c=1-q_c^{(0)}=0.40381381,
\]
the absorbing sink is of order
\[
-\lambda_{\rm esc}s_c\simeq-0.1657\,\mathcal K_0.
\]

The importance-sampled killed hard residual at the fold is
\[
H_{\rm hard}^{\rm kill}
=
-0.24318,
\]
with local temperature derivative
\[
\partial_\Theta H_{\rm hard}^{\rm kill}
=
+0.04805.
\]
Adding the same matched BPS reconstructed-FP finite term gives
\[
\boxed{
H_{\rm BPS}^{\rm kill}
=
-0.19799,
}
\]
\[
\boxed{
\partial_\Theta H_{\rm BPS}^{\rm kill}
=
+0.01914.
}
\]

Consequently the killed-population fold is
\[
\boxed{
q_c^{\rm kill}(L_{\rm C})
=
0.59618619
+
\frac{0.0683356}{L_{\rm C}}
+
O(L_{\rm C}^{-2}),
}
\]
\[
\boxed{
\Theta_c^{\rm kill}(L_{\rm C})
=
0.32464127
-
\frac{0.474482}{L_{\rm C}}
+
O(L_{\rm C}^{-2}),
}
\]
and
\[
\boxed{
e_c^{\rm kill}(L_{\rm C})
=
0.80284620
-
\frac{0.0507454}{L_{\rm C}}
+
O(L_{\rm C}^{-2}).
}
\]

Thus absorption produces a parametrically larger NLO displacement than the
bound-survivor hard correction.  It moves the endpoint toward larger $q$ and
lower $\Theta$.

If one shifts the endpoint to NLO but retains the leading-log curvature,
the local killed-moment zero set has the approximation
\[
F_{\rm kill}
\simeq
F_q\left(q-q_c^{\rm kill}\right)
+
\frac12F_{\Theta\Theta}
\left(\Theta-\Theta_c^{\rm kill}\right)^2.
\]
Therefore, on the side where the two branches exist,
\[
\boxed{
\Theta_\pm^{\rm kill}(q;L_{\rm C})
=
\Theta_c^{\rm kill}(L_{\rm C})
\pm
3.1385\,
\sqrt{q_c^{\rm kill}(L_{\rm C})-q}
+\cdots.
}
\]
The $3.1385$ coefficient is the LL curvature
$\sqrt{2F_q/F_{\Theta\Theta}}$; its own $O(L_{\rm C}^{-1})$ renormalization
would require $q$-derivatives of the full killed NLO function and is not
included here.

For illustration, if the BPS coupling convention is parameterized by
$L_{\rm C}=-\ln g$,
\[
\begin{array}{c|ccc}
g&q_c^{\rm kill}&\Theta_c^{\rm kill}&e_c^{\rm kill}\\
\hline
10^{-2}&0.611025&0.221609&0.791827\\
10^{-3}&0.606079&0.255953&0.795500\\
10^{-4}&0.603606&0.273125&0.797337\\
10^{-6}&0.601132&0.290297&0.799173
\end{array}
\]
The $g=10^{-2}$ row has $1/L_{\rm C}\simeq0.22$ and should be regarded as
qualitative rather than a precision asymptotic prediction.

This boundary is the zero set of the generator acting on the
\emph{cemetery-extended, unnormalized} observable $s$, with $s=0$ after
escape.  It therefore combines shape drift with probability loss.  It is not
the phase boundary of the normalized surviving population.  For a survivor
density normalized at each time, the evolution contains the nonlinear
selection term
\[
-[\lambda_{\rm esc}-\langle\lambda_{\rm esc}\rangle]\rho_{\rm surv},
\]
so a genuine conditional-population boundary requires an ensemble
calculation.  The killed-moment endpoint should therefore be reported
separately from both the local bound-survivor fold and any future
conditional-survivor phase diagram.

\section{Three-dimensional numerical status}

The three-dimensional leading-log fold was independently recomputed from the
local generator of $s(h,L)$ and Gauss--Legendre Kepler quadrature.  The result
\[
(q_c,\Theta_c)=(0.59618619,0.32464127)
\]
and the derivatives
\[
F_q=2.89734555,\qquad
F_{q\Theta}=3.80451347,\qquad
F_{\Theta\Theta}=0.58826836
\]
are stable under quadrature refinement and finite-difference changes.

The exact hard residual was evaluated with scrambled Sobol importance
sampling, exact Maxwellian bath sampling, uniform mean anomaly for the Kepler
time measure, and the grazing coordinate
$t=\sin^2(\theta/2)=y^2$.  Around the physical $s$ fold the quoted
$H_{\rm hard}$ values use four $2^{19}$-sample replicates.  The value of
$H_{\rm hard}$ is substantially better conditioned than its numerical
$\Theta$ derivative.

The BPS reconstructed-FP numbers $H_{\rm FP}$ and
$\partial_\Theta H_{\rm FP}$ come from the certified classical equal-mass OCP
BPS coefficient implementation with cancellation of the arbitrary matching
scale checked numerically.  A publication version should include the
coefficient-evaluation code or a compact reproducibility appendix; the
present draft records the numerical outputs but does not yet constitute an
interval-certified computation.

\section{Revised EFT hierarchy}

The audited theory now separates three levels:
\[
\boxed{
\begin{array}{c}
\text{planar oriented cone sector: }R_-^2\\
\downarrow\\
\text{3D isotropic LL shape sector: }s=1-q\\
\downarrow\\
\text{BPS NLO matching + exact hard jumps}\\
\downarrow\\
\text{bound-survivor fold and cemetery-weighted killed-moment endpoint}.
\end{array}
}
\]
The planar analysis remains useful as a constrained stochastic model and as
the historical route by which the sign-changing cone Wilson function was
found.  The physical isotropic-plasma completion, however, is the
three-dimensional bounded-shape theory.


\section{Conclusion}

The audit resolves three distinct layers of the bath-coupled Kepler problem.
First, correcting the half-covariance It\^o normalization overturns the old
planar monotonicity conjecture and produces a sign-changing $R_-^2$ response.
That planar sector is internally consistent and retains its extreme-
eccentricity radial-edge asymptotics, but it describes noise constrained to
the orbital plane.

Second, a physical isotropic three-dimensional bath requires the angular-
momentum magnitude $L=|\bm r\times\bm v|$.  The It\^o curvature of $L$
changes the leading-log generator, and exact hard collisions force a bounded
observable.  The resulting normalization chain admits an external analytic
benchmark.  As $\Theta\to0$,
\[
\boxed{
\left\langle\mathcal Ls\right\rangle
=
\frac{5D_0T}{2\mu q}(3q^2-1),
\qquad
q_0=\frac1{\sqrt3},
}
\]
which agrees identically with the high-speed isotropic-encounter result of
Hamilton \& Modak after matching their kick variance to
$\langle|\Delta\bm v|^2\rangle/\Delta t=6D_0$.  Thus the physical
finite-temperature result is anchored to an independent analytic limit rather
than only to internal numerical checks.  For the jump-safe shape $s=1-q$,
the continuation of this hot-limit zero terminates at the leading-log fold
\[
\boxed{
(q_c^{(0)},\Theta_c^{(0)},e_c^{(0)})
=
(0.59618619,\;0.32464127,\;0.80284620).
}
\]

Third, controlled plasma NLO matching is not a scalar Coulomb-log replacement.
The BPS reconstructed-FP finite term must be combined with the exact
short-distance hard-jump residual.  For the bound-preserving shape generator,
the hard residual reverses the direction of the formal FP shift and leaves
\[
\boxed{
q_c^{\rm bound}
=
0.59618619+\frac{0.0111477}{L_{\rm C}}+\cdots,
\qquad
\Theta_c^{\rm bound}
=
0.32464127+\frac{0.12838}{L_{\rm C}}+\cdots.
}
\]
For the cemetery-extended killed observable, escape produces the larger
unconditional moment shift
\[
\boxed{
q_c^{\rm kill}
=
0.59618619+\frac{0.0683356}{L_{\rm C}}+\cdots,
\qquad
\Theta_c^{\rm kill}
=
0.32464127-\frac{0.474482}{L_{\rm C}}+\cdots.
}
\]
The latter is not yet the phase boundary of the normalized surviving
population; that requires the hazard-selection term at the ensemble level.

The resulting open EFT is therefore not a single diffusion on a cone.  It is
a hierarchy of planar cone diffusion, three-dimensional jump-safe shape
dynamics, BPS-matched finite hard jumps, and an absorbing unbound channel.
That separation is the central structural result of the present draft.

\begin{thebibliography}{99}

\bibitem{Rosenbluth1957}
M.~N. Rosenbluth, W.~M. MacDonald, and D.~L. Judd,
``Fokker--Planck Equation for an Inverse-Square Force,''
\emph{Physical Review} \textbf{107}, 1--6 (1957).
doi:10.1103/PhysRev.107.1.

\bibitem{CohnKulsrud1978}
H. Cohn and R.~M. Kulsrud,
``The Stellar Distribution Around a Black Hole: Numerical Integration of the
Fokker--Planck Equation,''
\emph{The Astrophysical Journal} \textbf{226}, 1087--1108 (1978).

\bibitem{Lenard1960}
A. Lenard,
``On Bogoliubov's Kinetic Equation for a Spatially Homogeneous Plasma,''
\emph{Annals of Physics} \textbf{10}, 390--400 (1960).
doi:10.1016/0003-4916(60)90003-8.

\bibitem{Balescu1960}
R. Balescu,
``Irreversible Processes in Ionized Gases,''
\emph{Physics of Fluids} \textbf{3}, 52--63 (1960).
doi:10.1063/1.1706002.

\bibitem{BPS2005}
L.~S. Brown, D.~L. Preston, and R.~L. Singleton Jr.,
``Charged Particle Motion in a Highly Ionized Plasma,''
\emph{Physics Reports} \textbf{410}, 237--333 (2005).
doi:10.1016/j.physrep.2005.01.001.

\bibitem{CollinsSari2008}
B.~F. Collins and R. Sari,
``L\'evy Flights of Binary Orbits due to Impulsive Encounters,''
\emph{The Astronomical Journal} \textbf{136}, 2552--2562 (2008).
doi:10.1088/0004-6256/136/6/2552.

\bibitem{FouvryPichonMagorrian2017}
J.-B. Fouvry, C. Pichon, and J. Magorrian,
``The Secular Evolution of Discrete Quasi-Keplerian Systems. I. Kinetic
Theory of Stellar Clusters near Black Holes,''
\emph{Astronomy \& Astrophysics} \textbf{598}, A71 (2017).
doi:10.1051/0004-6361/201629645.

\bibitem{HamiltonModak2024}
C. Hamilton and S. Modak,
``Eccentricity Dynamics of Wide Binaries -- II. The Effect of Stellar
Encounters and Constraints on Formation Channels,''
\emph{Monthly Notices of the Royal Astronomical Society}
\textbf{532}, 2425--2440 (2024).
doi:10.1093/mnras/stae1654.

\end{thebibliography}

\end{document}
